\documentclass[11pt]{article}
\usepackage[T1]{fontenc}
\usepackage{textcomp}
\usepackage[margin=0.75in]{geometry}
\usepackage{amsmath,amssymb,amsthm}
\usepackage{array,booktabs}
\usepackage{hyperref}
\setlength{\parindent}{0pt}
\newtheorem{theorem}{Theorem}
\theoremstyle{definition}
\newtheorem{definition}{Definition}
\title{Flow Proof Artifact: prop-13-gnomon-from-two-segments}
\author{Generated by \texttt{flow doc proof}}
\date{Flow Proof Book}
\begin{document}
\maketitle
To apply a gnomon equal to the rectangle contained by two unequal straight lines.\\[0.5em]
\textbf{Source.} euclid — Elements, Book II, Proposition 13\\[0.5em]
\bigskip
\noindent{\large \textbf{Derived fact 1.} \textit{To apply a gnomon equal to the rectangle contained by two unequal straight lines} \hfill the Euclidean plane \cdot Euclid Book II \cdot Proposition 13: a gnomon equals the rectangle by two unequal segments}\par
\smallskip\noindent\textit{Coordinate: the Euclidean plane \cdot Euclid Book II \cdot Proposition 13: a gnomon equals the rectangle by two unequal segments}\par
\smallskip\noindent\textit{Source: euclid — Elements, Book II, Proposition 13}\par
\smallskip\noindent\textit{Built on: proposition 11: a gnomon equals the rectangle by the segments, for Euclid Book II on the Euclidean plane, proposition 5: the rectangle by unequal parts plus the square on the midpoint equals the square on the half, for Euclid Book II on the Euclidean plane}\par
\medskip
\noindent\textbf{Goal.} To apply a gnomon equal to the rectangle contained by two unequal straight lines\par
\noindent$\displaystyle \text{gnomon equals rect of unequal segments}$\par
\medskip
\noindent
\begin{tabular}{@{} >{\raggedright\arraybackslash}p{0.06\textwidth} >{\raggedright\arraybackslash}p{0.47\textwidth} >{\raggedleft\arraybackslash}p{0.41\textwidth} @{}}
\textbf{\#} & \textbf{Proof} & \textbf{Mathematics} \\
\midrule
\textbf{1.} & We prove that proposition 13: a gnomon equals the rectangle by two unequal segments for Euclid Book II the Euclidean plane. & \\[0.45em]
\textbf{2.} & Let unequal = segments\_AC\_and\_CB. & \\[0.45em]
\textbf{3.} & Let gnomon = gnomon\_about\_diameter. & \\[0.45em]
\textbf{4.} & From step 2 and step 3, we can deduce that gnomon equals rect(unequal). & $\displaystyle gnomon = rect(unequal)$ \\[0.45em]
\textbf{5.} & We invoke the derived fact governing Euclid Book II the Euclidean plane: proposition 11: a gnomon equals the rectangle by the segments, for Euclid Book II on the Euclidean plane. & \\[0.45em]
\textbf{6.} & We invoke the derived fact governing Euclid Book II the Euclidean plane: proposition 5: the rectangle by unequal parts plus the square on the midpoint equals the square on the half, for Euclid Book II on the Euclidean plane. & \\[0.45em]
\textbf{7.} & From step 5 and step 6, this implies gnomon equals rect of unequal segments. Hence proven. & $\displaystyle \text{gnomon equals rect of unequal segments}$ \\[0.45em]
\bottomrule
\end{tabular}
\label{thm:Geometry-Euclid Book II-proposition_13_a_gnomon_equals_the_rectangle_by_two_unequal_segments}
\par\medskip\hrule
\end{document}
